% !TEX root = ../paper.tex
\doclearpage
\section{Acceleration via Preprocessing}
\label{sec:preprocessing}

The protocols of \cref{sec:piop-for-pattern-matching,sec:lex-order} are complete as they stand, but both pay for their generality: pattern matching sweeps the character-level columns of every candidate string, and sorting inspects the strings character by character. In both cases much of this work can be shifted from the prover, who would repeat it for every query, to the data owner, who performs it once at ingestion and commits the result alongside the data. This section presents the two preprocessing steps we use. For pattern matching, \emph{fingerprints} (Sections~\ref{sec:prefilter-fp} to~\ref{sec:opt-fp}) summarize which characters each string contains, so that a cheap pre-filter discards the rows that cannot match and the full protocol runs only on the survivors. For ordering, \emph{ranks} (\cref{sec:rank}) record the position of every committed string in the lexicographic order of its column, so that sorting strings reduces to sorting integers.

The two steps differ in weight, and the difference is instructive. Ranks are a single extra column and need no new protocol, because the order of a column is a property of the column alone. Matching depends on the pattern, so its summary must serve every pattern at once, which is why fingerprints come with a design space, a PIOP of their own, and a cost model to decide when to use them. What the two share is the trust model: the precomputed columns are committed by the data owner, who is trusted for the data itself, and where this trust is unwanted they can instead be certified once at ingestion against the committed character-level columns, as \cref{piop:rank-cert} does for ranks, so that no query ever pays for their correctness.

% !TEX root = ../paper.tex
\subsection{Pre-Filtering via Fingerprinting}
\label{sec:prefilter-fp}

In \cref{sec:piop-for-pattern-matching} we presented a PIOP for pattern matching that is proven secure and ready to be deployed. Although \cref{piop:multi-char-pattern-matching} proves what is conceptually a simple filter, it operates on the full char-level machinery of the string arithmetization and is therefore expensive. This motivates us to shrink the string columns fed into it as much as possible.

To this end, we borrow the \emph{fingerprinting} technique from the database literature, where it was introduced~\cite{string_fingerprint} as a fast pre-filter for query execution: it cheaply discards \emph{true negatives}, the rows that certainly do not match the pattern, while possibly retaining \emph{false positives}, rows that survive the filter but do not actually match. Once most true negatives are gone, the surviving column is small enough that an exact but expensive pattern-matching protocol can affordably remove the remaining false positives.

Although fingerprinting was originally proposed to speed up the \emph{execution} of pattern matching, we find that it speeds up its \emph{verification} just as effectively. Before presenting the verification mechanism, we first describe how fingerprint-based pre-filtering works.

The idea of fingerprinting is to replace the full representation of a string with a handful of \emph{bins}, each bin holding a single bit that is $1$ if and only if the string satisfies a fixed feature. The fingerprint of a string is the vector of its bin values.

For example, suppose we have $m=4$ bins, where bin $i\in\{1,2,3,4\}$ is computed by its own predicate $\mathsf{fp}_i: \mathsf{Strings} \to \{0,1\}$ defined as follows:
\ifsp
%% S&P (two-column): one predicate per row so the display fits one column.
\begin{align*}
 \mathsf{fp}_1(\mathsf{s}) &= \begin{cases}
  1 & \text{if $\mathsf{s}$ contains `a', `b', or `c'}\\
  0 & \text{otherwise}
 \end{cases}\\
 \mathsf{fp}_2(\mathsf{s}) &= \begin{cases}
  1 & \text{if $\mathsf{s}$ contains any of `d', `e', \dots, `l'}\\
  0 & \text{otherwise}
 \end{cases}\\
 \mathsf{fp}_3(\mathsf{s}) &= \begin{cases}
  1 & \text{if $\mathsf{s}$ contains `m' or `n'}\\
  0 & \text{otherwise}
 \end{cases}\\
 \mathsf{fp}_4(\mathsf{s}) &= \begin{cases}
  1 & \text{if $\mathsf{s}$ contains any of `o', `p', \dots, `z'}\\
  0 & \text{otherwise}
 \end{cases}
\end{align*}
\else
\begin{align*}
 &\mathsf{fp}_1(\mathsf{s}) = \begin{cases}
  1 & \text{if $\mathsf{s}$ contains `a', `b', or `c'}\\
  0 & \text{otherwise}
 \end{cases}
 \qquad
 &\mathsf{fp}_2(\mathsf{s}) = \begin{cases}
  1 & \text{if $\mathsf{s}$ contains any of `d', `e', \dots, `l'}\\
  0 & \text{otherwise}
 \end{cases}
 \\
 &\mathsf{fp}_3(\mathsf{s}) = \begin{cases}
  1 & \text{if $\mathsf{s}$ contains `m' or `n'}\\
  0 & \text{otherwise}
 \end{cases}
 \qquad
 &\mathsf{fp}_4(\mathsf{s}) = \begin{cases}
  1 & \text{if $\mathsf{s}$ contains any of `o', `p', \dots, `z'}\\
  0 & \text{otherwise}
 \end{cases}
\end{align*}
\fi

Now suppose we want to evaluate the filter \texttt{\%dar\%} on the strings \texttt{darling}, \texttt{zebra}, and \texttt{lemon}. Fingerprinting the pattern's literal \texttt{dar} and the three strings yields
\ifsp
%% S&P (two-column): two fingerprints per row so the display fits one column.
\begin{align*}
 \texttt{dar} &\to 1101, & \texttt{darling} &\to 1111,\\
 \texttt{zebra} &\to 1101, & \texttt{lemon} &\to 0111.
\end{align*}
\else
\begin{equation*}
 \texttt{dar}\to 1101,\qquad \texttt{darling} \to 1111,\qquad \texttt{zebra} \to 1101,\qquad \texttt{lemon}\to 0111.
\end{equation*}
\fi

A string can match \texttt{\%dar\%} only if it contains every character of \texttt{dar}, so every bin that is active in the pattern's fingerprint must also be active in the string's fingerprint. The first bin is active in the pattern but inactive in \texttt{lemon}: \texttt{lemon} contains none of `a', `b', `c', so it cannot contain \texttt{dar} and is discarded immediately. Both \texttt{darling} and \texttt{zebra}, on the other hand, activate every bin that the pattern does, so both survive the filter. Of these, \texttt{darling} is a genuine match, whereas \texttt{zebra} does not satisfy the predicate and is a \emph{false positive} that only the exact protocol can remove.

\subsection{Fingerprint Design}
\label{sec:fp-design}
The effectiveness of pre-filtering hinges on the features assigned to each bin, so a crucial question is how to choose them. Prior works (see \cref{sec:fp-prior}) mostly aim to minimize the execution time of the fingerprinting-based filtering; However, our objective is to minimze the prover effort to certify ther this process. Therefore, we design our own fingerprinting as follows: Given a string column $\col$ as the corpus, the data owner extracts (1) the single-character, (2) bigram, (3) trigram, and (4) precedence features occurring in the corpus, where a string has the precedence feature $(a, b)$ for alphanumerics $a$ and $b$ if some occurrence of $a$ lies before the last occurrence of $b$, with the followin two restrictions:
\begin{itemize}[nosep]
  \item  A feature present in every row does not help us in filtering task, hence discarded;
\item If more than $m_{\max}$ features remain, the data owner keeps the $m_{\max}$ that occur in the most rows ($m_{\max} = 2048$ in our implementation).
\end{itemize}

% Bigrams and trigrams are read over an alphabet that keeps the alphanumerics and the space and merges every other character into a single symbol, so that an n-gram can also record where a word starts or ends, as patterns such as \texttt{\% i\%} require.
Every feature receives its own bin, i.e.\ $\mathsf{fp}_i(\mathsf{s}) = 1$ if and only if the string $\mathsf{s}$ has the $i$-th feature. This yields the fingerprint $\mathcal{F}_\col=\{\mathsf{fp_1},\dots,\mathsf{fp}_m\}$ where $m\leq m_{\max}$. 

Note that no optimization takes place at fingerprinting time; the data owner only extracts features. Instead, all bins are made available to the prover, who selects the most useful subset for the workload at hand, namely the pattern given at query time.

For each bin $i$, the data owner forms the binary \emph{bin column} $\mathsf{bin}_i = \bigl(\mathsf{fp}_i(\col[j])\bigr)_{j=1}^{N}$ and commits to its multilinear extension, the \emph{bin polynomial} $\MLE{\mathsf{bin}}_i$, separately; instead of publishing the $m$ commitments it publishes only the root of a Merkle tree whose $i$-th leaf is the commitment to $\MLE{\mathsf{bin}}_i$ (\cref{alg:fp-commit}). The published commitment therefore grows by a single hash per column, however large $m$ is. 

\begin{algorithm}[h]
\caption{Fingerprint construction by the data owner.}
\label{alg:fp-commit}
\begin{algorithmic}[1]
\Require string column $\col$ with $N$ rows; feature categories $\mathcal{K}$ (here $\mathcal{K} = \{\textsc{char}, \textsc{bigram}, \textsc{trigram}, \textsc{prec}\}$); bin budget $m_{\max}$
\Ensure the fingerprint $\mathcal{F}_\col = \{\mathsf{fp}_1, \dots, \mathsf{fp}_m\}$, the commitments $\mathsf{com}_1, \dots, \mathsf{com}_m$ to the bin polynomials $\MLE{\mathsf{bin}}_1, \dots, \MLE{\mathsf{bin}}_m$, and their Merkle root $\rho_\col$
\State $\mathcal{F}_\col \gets \emptyset$;\quad $\mathrm{cnt}[\cdot] \gets 0$
% \Comment{features seen, and rows having each}
\For{$i = 1, \dots, N$}
  \For{$f \in \mathrm{features}_{\mathcal{K}}(\col[i])$}
  % \Comment{each distinct feature of the row, once}
    \State $\mathcal{F}_\col \gets \mathcal{F}_\col \cup \{f\}$;\quad $\mathrm{cnt}[f] \gets \mathrm{cnt}[f] + 1$
  \EndFor
\EndFor
\State $\mathcal{F}_\col \gets \{f \in \mathcal{F}_\col : \mathrm{cnt}[f] < N\}$
% \Comment{a feature in every row discards no row}
\State sort $\mathcal{F}_\col$ by $\mathrm{cnt}$, descending
% \Comment{ties broken by a fixed order on features}
\State $m \gets \min(|\mathcal{F}_\col|, m_{\max})$;\quad $(f_1, \dots, f_m) \gets$ the first $m$ features of $\mathcal{F}_\col$
\For{$b = 1, \dots, m$} \Comment{one binary column per bin}
  \State $\mathsf{fp}_b(\mathsf{s}) := [\,f_b \in \mathrm{features}_{\mathcal{K}}(\mathsf{s})\,]$ \Comment{bin predicate}
  \State $\mathsf{bin}_b \gets \bigl(\mathsf{fp}_b(\col[i])\bigr)_{i=1}^{N}$;\quad $\mathsf{com}_b \gets \mathsf{Commit}(\MLE{\mathsf{bin}}_b)$
\EndFor
\State $\rho_\col \gets \mathsf{MerkleRoot}(\mathsf{com}_1, \dots, \mathsf{com}_m)$
\State \Return $\mathcal{F}_\col \gets \{\mathsf{fp}_1, \dots, \mathsf{fp}_m\}$, $(\mathsf{com}_1, \dots, \mathsf{com}_m)$ and $\rho_\col$
\end{algorithmic}
\end{algorithm}

\subsection{Fingerprinting PIOP}
\label{sec:fp-piop}




Let $\col$ be a string column with $N$ rows whose bins $\mathsf{fp}_1,\dots,\mathsf{fp}_m$ the data owner committed under the Merkle root $\rho_\col$ (\cref{alg:fp-commit}), and let $\mathcal{F}_\zeta=\{b : \mathsf{fp}_b(\zeta)=1\}$ be the bins of pattern $\zeta$. The Pre-Filtering Check (\cref{piop:prefilter-check}) proves membership in the relation $\mathcal{R}_{\mathsf{PF}}$, stated in \cref{sec:proof-pf}: the prover chooses a non-empty subset $\mathcal{F}_u \subseteq \mathcal{F}_\zeta$ of the pattern\textquotesingle s bins and opens their commitments under $\rho_\col$; the new str-level activator is the old one multiplied by the opened bin columns, $\mathsf{a}^\prime[i] = \mathsf{a}[i]\cdot\prod_{b\in\mathcal{F}_u}\mathsf{fp}_b[i]$; and the new char-level activator keeps exactly the characters of the kept strings. Since every matching string lies in every bin of the pattern, the filter never discards a matching row (\cref{lem:prefilter-sound}).


\begin{DefinePIOP}{Pre-Filtering Check}{prefilter-check}[Pre-Filtering Check]
  \emergencystretch=3em\relax
  \emph{Inputs:} the pattern $\zeta$ and its bins $\mathcal{F}_\zeta=\{b : \mathsf{fp}_b(\zeta)=1\}$, which both parties derive from the public pattern; the Merkle root $\rho_\col$ over the bin commitments of column $\col$ (\cref{alg:fp-commit}); the input activators $\Oracle{\MLE{\mathsf{a}}}$ and $\Oracle{\MLE{\mathsf{char\text{-}act}}}$; and the fixed columns $\Oracle{\MLE{\mathsf{orig\text{-}ind}}}, \Oracle{\MLE{\mathsf{ind}}}, \Oracle{\MLE{\mathsf{l}}}$.
  \vspace{2mm}
  \begin{enumerate}[nosep]
    \item (\emph{Bin selection})\label[step]{pf:select} $\PIOPProver$ chooses the bins to test, $\mathcal{F}_u$, via \cref{alg:greedy-bins} and sends them to $\PIOPVerifier$. $\PIOPVerifier$ checks that $\mathcal{F}_u\neq\emptyset$ and $\mathcal{F}_u\subseteq\mathcal{F}_\zeta$, and rejects otherwise.
    \item (\emph{Opening the bins})\label[step]{pf:open} $\PIOPProver$ sends the commitments $(\mathsf{com}_b)_{b\in\mathcal{F}_u}$ together with a Merkle multiproof $\pi$ for the leaves $\mathcal{F}_u$. $\PIOPVerifier$ checks that $\pi$ opens $\rho_\col$ to $\mathsf{com}_b$ at leaf $b$ for every $b\in\mathcal{F}_u$, and rejects otherwise. From then on the parties hold the bin oracles $\Oracle{\MLE{\mathsf{fp}_b}}$ for $b\in\mathcal{F}_u$.
    \item (\emph{Filtered activators})\label[step]{pf:filter} $\PIOPProver$ computes the str-level activator $\mathsf{a'}[i] := \mathsf{a}[i]\cdot\prod_{b\in\mathcal{F}_u}\mathsf{fp}_b[i]$. It also computes its char-level counterpart $\mathsf{char\text{-}act'}$, which keeps exactly the active characters of the strings kept by $\mathsf{a'}$. $\PIOPProver$ sends $\Oracle{\MLE{\mathsf{a'}}}$ and $\Oracle{\MLE{\mathsf{char\text{-}act'}}}$ to $\PIOPVerifier$.
    \item (\emph{Subset test})\label[step]{pf:subset} The parties invoke a \cref{piop:zerocheck}, of degree $|\mathcal{F}_u|+1$, on $\mathsf{a'}-\mathsf{a}\cdot\prod_{b\in\mathcal{F}_u}\mathsf{fp}_b$.
    \item (\emph{Activator update validity})\label[step]{pf:dpuc} The parties invoke \cref{piop:data-preserving-update-check} on the updated column $\strcol'$, binding $\mathsf{char\text{-}act'}$ to $\mathsf{a'}$.
    \item (\emph{Output})\label[step]{pf:output} $\PIOPProver$ and $\PIOPVerifier$ output the filtered activators $\mathsf{a'}$ and $\mathsf{char\text{-}act'}$.
  \end{enumerate}
\end{DefinePIOP}

\subsection{Optimum usage of fingerprints}
\label{sec:opt-fp}

As mentioned in \cref{sec:fp-design}, all possible single-char, bigram, trigram, and proceeding features are extracted at the time of commitment and handed to the prover and verifier. However, we observe that it is not necessary to use all of these bins. To demonstrate this, we present two examples (1) consider the pattern $\zeta = \texttt{\%dar\%}$ over a column that contains the string \texttt{zebra}, so the data owner has committed a bin $\mathsf{fp}_{\texttt{z}}$ for the single-character feature \texttt{z}. The pattern, however, contains no \texttt{z} which means using the fingerprint would not remove any true negatives. (2) On the other hand, even if a fingerprint does exist in a pattern, adding it might have no effect. For example, both the trigram \texttt{dar} and the character \texttt{d} are features of $\zeta = \texttt{\%dar\%}$, and consider the column $(\texttt{darling}, \texttt{dare}, \texttt{dog}, \texttt{lemon}, \texttt{sun})$. Each bin removes true negatives on its own: $\mathsf{fp}_{\texttt{dar}}$ alone discards \texttt{dog}, \texttt{lemon} and \texttt{sun}, and $\mathsf{fp}_{\texttt{d}}$ alone discards \texttt{lemon} and \texttt{sun}. Yet once the column has been filtered by $\mathsf{fp}_{\texttt{dar}}$, only \texttt{darling} and \texttt{dare} survive, and applying $\mathsf{fp}_{\texttt{d}}$ afterwards discards nothing further, since every string containing \texttt{dar} also contains \texttt{d}.

This becomes more important when we notice that according to \cref{piop:prefilter-check}, each added bin adds prover, verifier, and proof size overheads too. The prover has to send the opening of the leaves of the Merkle Tree to the verifier. Moreover, the prover has to prove that the fingerprint filter was done correctly. That's why the prover needs to find out the minimal set of bins $\mathcal{F}_u$ that minimizes the prover time. anything beyound that won't increase the selectivity, hence won't help with the proving time, yet increases other types of overheads.


Finding such a set on the prover side in runtime, comes with two barriers:
\begin{itemize}[nosep]
  \item The prover cannot test the prover times online, i.e. it cannot afford to run the prover multiple times and pick the best one. This defeats the whole purpose of minimizing the prover time 
  \item Finding the best possible bins requires trying out each possible set of bins in $\mathcal{F}_\zeta$ which takes $2^{|\mathcal{F}_\zeta|}$ and is potentially large.
\end{itemize}
We fix the first issue in \cref{sec:cost-model} by using a \emph{cost model}, instead of running the actual prover, and we fix the second issue in \cref{sec:greedy} by using a \emph{greedy} assignment for placing the bins and empirically show that the result is very close to the best possible fingerprint.


\subsubsection{Prover Cost Model}
\label{sec:cost-model}
To prevent the prover to actually run the prover each time it wants to test a new set of bins, we model the cost of string filtering from its \emph{shape} as follow.

\parhead{Notation} The input rows are denoted by $N$ and characters by $M$, the surviving rows by $S$ and their characters by $M_S$, the number of matching rows $T$, the number of tested bins $n_{\mathrm{bin}} = |\mathcal{F}_u|$, the pattern's total literal length $L_\zeta$, and its factor count $n_\zeta$. Write $D(x)$ for the smallest power of two at least $x$, and $D^+(x)$ for the strict variant, which doubles an exact power of two. Tables are padded to a power of two, so the unfiltered input occupies the domain $N_{\mathrm{pad}} = D(N)$. \Cref{fig:plan-shape} places these parameters on the query plan.

\begin{figure}[t]
\centering
\resizebox{\ifsp\columnwidth\else\linewidth\fi}{!}{%
\begin{tikzpicture}[
  scell/.style={draw, minimum width=1.3cm, minimum height=0.42cm, inner sep=0pt, font=\ttfamily\scriptsize},
  ccell/.style={draw, minimum width=0.42cm, minimum height=0.26cm, inner sep=0pt, font=\ttfamily\tiny},
  bnd/.style={very thick},
  pad/.style={dashed, draw=gray},
  surv/.style={fill=orange!25},
  match/.style={fill=green!25},
  hdr/.style={font=\scriptsize},
  stage/.style={font=\small\bfseries, align=center},
  param/.style={font=\small, align=center},
  flow/.style={-{Stealth[length=2.5mm]}, thick}
]
% ---------------- input column ----------------
\node[stage] at (1.3,0.95) {Input column};
\node[hdr] at (0.7,0.42) {$\mathsf{str}$}; \node[hdr] at (1.95,0.42) {$\mathsf{char}$};
\node[scell,] at (0.7,-0) {fig};
\node[scell,surv] at (0.7,-0.42) {dart};
\node[scell,] at (0.7,-0.84) {sun};
\node[scell,surv] at (0.7,-1.26) {zebra};
\node[scell,surv] at (0.7,-1.68) {dare};
\node[ccell,] at (1.95,-0) {f};
\draw[bnd] (1.74,0.13) -- (2.16,0.13);
\node[ccell,] at (1.95,-0.26) {i};
\node[ccell,] at (1.95,-0.52) {g};
\node[ccell,surv] at (1.95,-0.78) {d};
\draw[bnd] (1.74,-0.65) -- (2.16,-0.65);
\node[ccell,surv] at (1.95,-1.04) {a};
\node[ccell,surv] at (1.95,-1.3) {r};
\node[ccell,surv] at (1.95,-1.56) {t};
\node[ccell,] at (1.95,-1.82) {s};
\draw[bnd] (1.74,-1.69) -- (2.16,-1.69);
\node[ccell,] at (1.95,-2.08) {u};
\node[ccell,] at (1.95,-2.34) {n};
\node[ccell,surv] at (1.95,-2.6) {z};
\draw[bnd] (1.74,-2.47) -- (2.16,-2.47);
\node[ccell,surv] at (1.95,-2.86) {e};
\node[ccell,surv] at (1.95,-3.12) {b};
\node[ccell,surv] at (1.95,-3.38) {r};
\node[ccell,surv] at (1.95,-3.64) {a};
\node[ccell,surv] at (1.95,-3.9) {d};
\draw[bnd] (1.74,-3.77) -- (2.16,-3.77);
\node[ccell,surv] at (1.95,-4.16) {a};
\node[ccell,surv] at (1.95,-4.42) {r};
\node[ccell,surv] at (1.95,-4.68) {e};
\node[param] at (0.7,-2.24) {$N$ rows};
\node[param] at (1.95,-5.16) {$M$ chars};
% ---------------- pre-filter ----------------
\draw[flow] (2.8,-0.63) -- (5.4,-0.63) node[midway,above,stage] {Pre-filter} node[midway,below,param] {tests $n_{\mathrm{bin}}$\\fingerprint bins};
% ---------------- rematerialized column ----------------
\node[stage] at (7.15,0.95) {Rematerialized \normalfont\small(if $\mathrm{remat}$)};
\node[hdr] at (6.4,0.42) {$\mathsf{str}$}; \node[hdr] at (7.9,0.42) {$\mathsf{char}$};
\node[scell,surv] at (6.4,-0) {dart};
\node[scell,surv] at (6.4,-0.42) {zebra};
\node[scell,surv] at (6.4,-0.84) {dare};
\node[scell,pad] at (6.4,-1.26) {};
\node[ccell,surv] at (7.9,-0) {d};
\draw[bnd] (7.69,0.13) -- (8.11,0.13);
\node[ccell,surv] at (7.9,-0.26) {a};
\node[ccell,surv] at (7.9,-0.52) {r};
\node[ccell,surv] at (7.9,-0.78) {t};
\node[ccell,surv] at (7.9,-1.04) {z};
\draw[bnd] (7.69,-0.91) -- (8.11,-0.91);
\node[ccell,surv] at (7.9,-1.3) {e};
\node[ccell,surv] at (7.9,-1.56) {b};
\node[ccell,surv] at (7.9,-1.82) {r};
\node[ccell,surv] at (7.9,-2.08) {a};
\node[ccell,surv] at (7.9,-2.34) {d};
\draw[bnd] (7.69,-2.21) -- (8.11,-2.21);
\node[ccell,surv] at (7.9,-2.6) {a};
\node[ccell,surv] at (7.9,-2.86) {r};
\node[ccell,surv] at (7.9,-3.12) {e};
\node[ccell,pad] at (7.9,-3.38) {};
\node[ccell,pad] at (7.9,-3.64) {};
\node[ccell,pad] at (7.9,-3.9) {};
\node[param] at (6.4,-2.02) {$S$ rows\\padded to $D(S)$};
\node[param] at (7.9,-4.58) {$M_S$ chars\\padded to $D(M_S)$};
\node[param] at (7.15,-5.38) {$(N', M') = (S, M_S)$};
% ---------------- pattern matching ----------------
\draw[flow] (8.6,-0.63) -- (10.6,-0.63) node[midway,above,stage] {Pattern\\matching} node[midway,below,param] {$\zeta = \texttt{\%dar\%}$\\$L_\zeta$, $n_\zeta$};
% ---------------- matches ----------------
\node[stage] at (12.15,0.95) {Matches};
\node[hdr] at (11.4,0.42) {$\mathsf{str}$}; \node[hdr] at (12.9,0.42) {$\mathsf{char}$};
\node[scell,match] at (11.4,-0) {dart};
\node[scell,match] at (11.4,-0.42) {dare};
\node[ccell,match] at (12.9,-0) {d};
\draw[bnd] (12.69,0.13) -- (13.11,0.13);
\node[ccell,match] at (12.9,-0.26) {a};
\node[ccell,match] at (12.9,-0.52) {r};
\node[ccell,match] at (12.9,-0.78) {t};
\node[ccell,match] at (12.9,-1.04) {d};
\draw[bnd] (12.69,-0.91) -- (13.11,-0.91);
\node[ccell,match] at (12.9,-1.3) {a};
\node[ccell,match] at (12.9,-1.56) {r};
\node[ccell,match] at (12.9,-1.82) {e};
\node[param] at (11.4,-0.98) {$T$ rows};
\node[param] at (12.9,-2.3) {$\approx M_T$ chars};
% ---------------- bypass (no rematerialization) ----------------
\draw[flow,dashed,gray] (1.3,1.3) -- (1.3,1.85) -- (12.15,1.85) -- (12.15,1.3);
\node[param,gray,fill=white,inner sep=2pt] at (7.0,1.85) {no $\mathrm{remat}$: pattern matching runs on $(N', M') = (N, M)$};
\end{tikzpicture}%
}
\caption{Shape of a pattern-matching plan with fingerprint pre-filtering, with the size parameters of the cost model. Each stage shows the str-level column $\mathsf{str}$ next to its character column $\mathsf{char}$, with thick rules marking string boundaries. The pre-filter tests $n_{\mathrm{bin}}$ fingerprint bins and keeps $S$ of the $N$ input rows (shaded). If $\mathrm{remat}$ holds, the survivors are compacted into new columns padded to $D(S)$ rows and $D(M_S)$ characters (dashed cells), and pattern matching runs on $(N', M') = (S, M_S)$; otherwise it runs on the full input (dashed path). Pattern matching then keeps the $T$ matching rows, discarding the survivors that were false positives.}
\label{fig:plan-shape}
\end{figure}

 All of these are known to the prover at planning time. The shape determines what the planner does:
\ifsp
%% S&P (two-column): one definition per row so the display fits one column.
\begin{align*}
  \mathrm{remat}  &= [n_{\mathrm{bin}} > 0] \wedge [D^+(S) < D^+(N_{\mathrm{pad}})],\\
  (N', M', P)     &= \begin{cases} (S, M_S, S) & \text{if } \mathrm{remat},\\
                                    (N, M, N_{\mathrm{pad}}) & \text{otherwise,} \end{cases}\\
  \mathrm{second} &= [D^+(T) < D^+(P)],\\
  M_T             &= T\, M' / N'.
\end{align*}
\else
\begin{align*}
  \mathrm{remat}  &= [n_{\mathrm{bin}} > 0] \wedge [D^+(S) < D^+(N_{\mathrm{pad}})], &
  (N', M', P)     &= \begin{cases} (S, M_S, S) & \text{if } \mathrm{remat},\\
                                    (N, M, N_{\mathrm{pad}}) & \text{otherwise,} \end{cases}\\
  \mathrm{second} &= [D^+(T) < D^+(P)], &
  M_T             &= T\, M' / N'.
\end{align*}
\fi
The survivors are rematerialized exactly when $\mathrm{remat}$ holds, that is, when a pre-filter is applied and the survivors fit in a strictly smaller power-of-two domain than the input. The pattern-matching protocol then runs on $N'$ rows and $M'$ characters. Every rematerialization of a string column carries a data-preserving update check (\cref{piop:data-preserving-update-check}) that binds the compacted characters to their source rows, so its cost scales with the \emph{character} domain of the survivors.

\parhead{The model} The predicted proving time is linear in the following terms, each priced on the padded domain that the corresponding stage actually runs on:
\ifsp
%% S&P (two-column): each description sits right after its term, shortened and
%% set slightly smaller, so the display fits one column.
\newcommand{\cmnote}[1]{\enspace\text{\small(#1)}}
\begin{align*}
  C = {}& a_0 \cmnote{fixed cost} \\
  &+ a_{\mathrm{pf}}\, [n_{\mathrm{bin}} > 0] \cmnote{pre-filter, if any bin tested} \\
  &+ c_{\mathrm{bin}}\, n_{\mathrm{bin}} \cmnote{per tested bin} \\
  &+ \mathrm{remat} \cdot a_r \cmnote{remat., fixed} \\
  &+ \mathrm{remat} \cdot a_{rs}\, D(S) \cmnote{remat., per row} \\
  &+ \mathrm{remat} \cdot a_{rc}\, D(M_S) \cmnote{consistency, per char} \\
  &+ a_c\, D(M') \cmnote{matching, per char} \\
  &+ a_{cL}\, (L_\zeta - 1)\, D(M') \cmnote{per char, per extra literal char} \\
  &+ a_{cn}\, (n_\zeta - 1)\, D(M') \cmnote{per char, per extra factor} \\
  &+ a_w\, D(N') \cmnote{matching, per row} \\
  &+ a_{wL}\, (L_\zeta - 1)\, D(N') \cmnote{per row, per extra literal char} \\
  &+ a_{wn}\, (n_\zeta - 1)\, D(N') \cmnote{per row, per extra factor} \\
  &+ \mathrm{second} \cdot \bigl(a_2 + a_{2r}\, D(N') \\
  &\phantom{{}+ \mathrm{second} \cdot \bigl(}{} + a_{2c}\, D(M_T)\bigr) \cmnote{small match set}
\end{align*}
\else
\begin{align*}
  C = {}& a_0
    && \text{(fixed cost of every proof)} \\
  &+ a_{\mathrm{pf}}\, [n_{\mathrm{bin}} > 0]
    && \text{(pre-filter check, paid once if any bin is tested)} \\
  &+ c_{\mathrm{bin}}\, n_{\mathrm{bin}}
    && \text{(opening and zerocheck factor per tested bin)} \\
  &+ \mathrm{remat} \cdot a_r
    && \text{(fixed cost of rematerializing the survivors)} \\
  &+ \mathrm{remat} \cdot a_{rs}\, D(S)
    && \text{(rematerialization, per surviving row)} \\
  &+ \mathrm{remat} \cdot a_{rc}\, D(M_S)
    && \text{(consistency check, per surviving character)} \\
  &+ a_c\, D(M')
    && \text{(pattern matching, per input character)} \\
  &+ a_{cL}\, (L_\zeta - 1)\, D(M')
    && \text{(per character, per extra literal character)} \\
  &+ a_{cn}\, (n_\zeta - 1)\, D(M')
    && \text{(per character, per extra factor)} \\
  &+ a_w\, D(N')
    && \text{(pattern matching, per input row)} \\
  &+ a_{wL}\, (L_\zeta - 1)\, D(N')
    && \text{(per row, per extra literal character)} \\
  &+ a_{wn}\, (n_\zeta - 1)\, D(N')
    && \text{(per row, per extra factor)} \\
  &+ \mathrm{second} \cdot \bigl(a_2 + a_{2r}\, D(N') + a_{2c}\, D(M_T)\bigr)
    && \text{(stage triggered by a small match set).}
\end{align*}
\fi


\subsubsection{Greedy Bin Assignment}
\label{sec:greedy}
The cost model prices a given set of bins, but the prover still has to decide which of the $2^{|\mathcal{F}_\zeta|}$ subsets of the pattern's bins to price. Since the prover holds the data, it knows for every bin $b$ the set $R_b$ of rows whose fingerprint has $\mathsf{fp}_b = 1$, and a set of bins $Q$ keeps exactly the rows in $\bigcap_{b \in Q} R_b$. The prover therefore orders the bins into a chain greedily (\cref{alg:greedy-bins}): starting from all rows, it repeatedly adds the bin that leaves the fewest survivors among the rows that are still alive, breaking ties toward the lower bin index. Choosing by survivors among the alive rows, rather than by each bin's own selectivity, avoids spending a bin on rows that earlier bins have already discarded: two bins that each discard many rows, but largely the same ones, are worth little more than one of them. After each step the prover records the query's shape, namely $S$, $M_S$ and $n_{\mathrm{bin}}$, which together with $N$, $M$, $T$, $L_\zeta$ and $n_\zeta$ is everything the cost model needs, and prices it. It then tests the prefix of the chain with the lowest predicted cost. This may be the empty prefix, in which case no pre-filter is applied; on a tie it takes the shorter prefix, since every tested bin also enlarges the proof.

The whole chain is priced rather than stopped at the first bin that does not pay off, because the cost is not convex in the number of bins: it drops in steps whenever $S$ or $M_S$ crosses a power of two, so a bin that saves nothing may precede one that crosses a step. Building the chain costs $O(|\mathcal{F}_\zeta|^2 N)$ bit operations on bitsets over the rows, which is negligible next to proving (about one second per query in our experiments). As discussed above, any prefix is sound, so the choice needs no agreement with the verifier beyond the list of bins, and an inaccurate cost model costs proving time, never correctness.

\begin{algorithm}[t]
\caption{Greedy bin assignment for pattern $\zeta$ over column $\col$.}
\label{alg:greedy-bins}
\begin{algorithmic}[1]
\Require pattern bins $\mathcal{F}_\zeta$; row sets $R_b = \{i : \mathsf{fp}_b(\col[i]) = 1\}$ for $b \in \mathcal{F}_\zeta$; cost model $C$
\Ensure bins to test $\mathcal{F}_u \subseteq \mathcal{F}_\zeta$
\State $A \gets \{1, \dots, N\}$ \Comment{rows still alive}
\State $Q_0 \gets \emptyset$;\quad $c_0 \gets C(S{=}N,\ M_S{=}M,\ n_{\mathrm{bin}}{=}0)$ \Comment{no pre-filter}
\For{$k = 1, \dots, |\mathcal{F}_\zeta|$}
  \State $b_k \gets \arg\min_{b \in \mathcal{F}_\zeta \setminus Q_{k-1}} |A \cap R_b|$ \Comment{ties to the lowest bin}
  \State $Q_k \gets Q_{k-1} \cup \{b_k\}$;\quad $A \gets A \cap R_{b_k}$
  \State $c_k \gets C\bigl(S{=}|A|,\ M_S{=}\textstyle\sum_{i \in A} |\col[i]|,\ n_{\mathrm{bin}}{=}k\bigr)$
\EndFor
\State $k^\star \gets \min\, \arg\min_{k} c_k$ \Comment{cheapest prefix, shortest on a tie}
\State \Return $\mathcal{F}_u \gets Q_{k^\star}$
\end{algorithmic}
\end{algorithm}

% !TEX root = ../paper.tex
\subsection{Precomputed Ranks for Ordering}
\label{sec:rank}

\Cref{piop:lex-sorted-check} inspects the characters of the sorted column, one position per round. In a relational plan, however, the strings being sorted are almost never new: every operator between a base table and the \texttt{ORDER BY} is data- or domain-preserving (\cref{sec:consistency-checking}), so each string in the sorted column is a copy of a string in a committed base column, and the \emph{relative order} of those strings is a property of the base column alone. We can therefore compute that order once, at ingestion, and store it next to each string. This is the rank column $\mathsf{rk}$ of \cref{sec:string-arith}: $\mathsf{rk}[i]$ is the dense rank of the $i$-th string among the distinct strings of its column, so that
\begin{displaymath}
  \begin{aligned}
    \mathsf{rk}[i] < \mathsf{rk}[j] &\iff s_i \lexlt s_j, \\
    \mathsf{rk}[i] = \mathsf{rk}[j] &\iff s_i = s_j .
  \end{aligned}
\end{displaymath}
Like the hash and the length, the rank is a str-level column inherited through every domain-preserving update by the provenance lookup of \cref{piop:domain-preserving-update-check}, so it is available, and trusted, wherever the string itself is. Two remarks on the definition. First, the rank must be \emph{dense}: if equal strings received distinct ranks, two rows with the same string would be forced into a fixed relative order instead of being ordered by the next sort key, which breaks the multi-key semantics of \cref{sec:sql-string-order}. Second, ranks are only comparable within one base column; the plan's lineage tells the planner which base column a sort key descends from, and two keys that descend from the same base column (such as the two copies of \texttt{n\_name} that Q7 of TPC-H joins and orders by) share a rank space.

\parhead{Rank validity}
A rank column is valid when it is monotone in the lexicographic order and separates exactly the distinct strings. Since $\lexle$ is a total preorder, it suffices to check this on the column \emph{sorted by rank}: writing $\pi$ for the permutation that sorts $\mathsf{rk}$ non-decreasingly, the relation $\NPRelation_{\mathsf{RankValid}}$ (stated in \cref{sec:proof-rankvalid}) requires that along $\pi$ the strings be non-decreasing in $\lexle$ and that two adjacent ranks be equal exactly when the two adjacent strings are equal. Adjacent consistency along $\pi$ implies the global property by transitivity: ranks are non-decreasing along the lexicographic order and increase exactly where the string changes.

\parhead{Certifying the ranks}
As discussed at the start of this section, the data owner may simply commit the rank column alongside the hash column and the fingerprint bins as part of the data. If instead the ranks are to be \emph{verified}, the char-level protocol of \cref{sec:lex-sorted-check} is exactly the tool: the prover materializes the base column sorted by rank, proves it is a permutation of the base column, proves the ranks non-decreasing along it, proves it lexicographically sorted, and finally proves that rank ties coincide with string ties, using the equal-neighbour indicator returned by \cref{piop:lex-sorted-check}, which marks precisely the adjacent pairs of equal strings. This is done once per committed string column, at ingestion, and amortizes over every query that later sorts a descendant of that column.

\begin{DefinePIOP}{Rank Certification Check}{rank-cert}[Rank Certification Check]
  \emergencystretch=3em\relax
  \emph{Inputs:} the committed column $\Oracle{\MLE{\strcol}}$ with its rank oracle $\Oracle{\MLE{\mathsf{rk}}}$; and $\Oracle{\MLE{\mathsf{ind}}}$.
  \vspace{2mm}
  \begin{enumerate}[nosep]
    \item (\emph{Sorted copy})\label[step]{rc:copy} $\PIOPProver$ materializes the column $\strcol^{\pi}$ obtained by sorting $\strcol$ by rank, with its rank column $\mathsf{rk}^{\pi}$, and sends all of its oracles. $\PIOPProver$ and $\PIOPVerifier$ invoke \cref{piop:domain-preserving-update-check} on $(\strcol, \strcol^{\pi})$, and \cref{piop:mset} on $(\mathsf{a}^{\pi}, \mathsf{h}^{\pi})$ and $(\mathsf{a}, \mathsf{h})$, as in \cref{ss:valid,ss:perm}.
    \item (\emph{Ranks non-decreasing})\label[step]{rc:mono} $\PIOPProver$ sends $\Oracle{\MLE{\mathsf{rk}^{\pi,\mathsf{next}}}}$ and the parties invoke \cref{piop:rotcheck} on $(\mathsf{rk}^{\pi}, \mathsf{rk}^{\pi,\mathsf{next}}, -1)$, then the $\mathsf{NonNeg}$ \cref{piop:sign-check} on the column with activator $\contig_{\mu,n-1}$ and data $\mathsf{rk}^{\pi,\mathsf{next}} - \mathsf{rk}^{\pi}$.
    \item (\emph{Strings sorted})\label[step]{rc:lex} $\PIOPProver$ and $\PIOPVerifier$ invoke \cref{piop:lex-sorted-check} on $\strcol^{\pi}$, obtaining its equal-neighbour indicator $\mathsf{same}^{\pi}$, which marks the adjacent pairs with $s^{\pi}_{i} = s^{\pi}_{i+1}$.
    \item (\emph{Ties coincide})\label[step]{rc:ties} $\PIOPProver$ and $\PIOPVerifier$ invoke \cref{piop:zerocheck} on $\mathsf{same}^{\pi} \mult (\mathsf{rk}^{\pi,\mathsf{next}} - \mathsf{rk}^{\pi})$, and the $\mathsf{Pos}$ \cref{piop:sign-check} on the column with activator $\contig_{\mu,n-1} \mult (1 - \mathsf{same}^{\pi})$ and data $\mathsf{rk}^{\pi,\mathsf{next}} - \mathsf{rk}^{\pi}$.
  \end{enumerate}
\end{DefinePIOP}

\Cref{rc:mono} is implied by \cref{rc:ties} and is listed only for clarity; it may be dropped. Density of the ranks (that they are exactly $0, \dots, D-1$ for $D$ distinct strings) is not needed for any protocol and is not checked; an honest data owner produces dense ranks, and any rank column in $\NPRelation_{\mathsf{RankValid}}$ orders the strings correctly.

\parhead{Sorting by rank}
With a trusted rank column in hand, the \texttt{ORDER BY} operator is proven exactly as in \cref{piop:str-sort}, except that the char-level \cref{ss:order} is replaced by an integer sortedness check on the inherited ranks: by $\NPRelation_{\mathsf{RankValid}}$, non-decreasing ranks along the output are equivalent to $s_i \lexle s_{i+1}$.

\begin{DefinePIOP}{Rank-Based String Sort Check}{rank-sort}[Rank-Based String Sort Check]
  \emergencystretch=3em\relax
  \emph{Inputs:} the input column $\Oracle{\MLE{\strcol^{\mathsf{old}}}}$ and the output column $\Oracle{\MLE{\strcol^{\mathsf{new}}}}$, both carrying rank oracles $\Oracle{\MLE{\mathsf{rk}^{\mathsf{old}}}}, \Oracle{\MLE{\mathsf{rk}^{\mathsf{new}}}}$; and $\Oracle{\MLE{\mathsf{ind}}}$.
  \vspace{2mm}
  \begin{enumerate}[nosep]
    \item (\emph{Validity, provenance, permutation})\label[step]{rs:valid} As in \cref{ss:valid,ss:perm}; the provenance lookup of \cref{piop:domain-preserving-update-check} transfers $\mathsf{rk}^{\mathsf{old}}$ to $\mathsf{rk}^{\mathsf{new}}$.
    \item (\emph{Order})\label[step]{rs:order} $\PIOPProver$ sends $\Oracle{\MLE{\mathsf{rk}^{\mathsf{new},\mathsf{next}}}}$ and the parties invoke \cref{piop:rotcheck} on $(\mathsf{rk}^{\mathsf{new}}, \mathsf{rk}^{\mathsf{new},\mathsf{next}}, -1)$, then the $\mathsf{NonNeg}$ \cref{piop:sign-check} on the column with activator $\contig_{\mu,n-1}$ and data $\mathsf{rk}^{\mathsf{new},\mathsf{next}} - \mathsf{rk}^{\mathsf{new}}$.
  \end{enumerate}
\end{DefinePIOP}

The online cost of the string sort is thus that of $\truthtable$'s integer sort, independent of the string lengths and of the char-level columns, which the planner may already have dropped (\cref{sec:string-arith}). Multi-key sorts and descending order carry over verbatim from \cref{sec:lex-sorted-check}, with the tie indicator $\mathbf{1}[\mathsf{rk}^{\mathsf{new},\mathsf{next}} = \mathsf{rk}^{\mathsf{new}}]$, since equal ranks are equal strings.

\parhead{Scope}
Ranks exist only for strings that descend from a committed column. A sort key produced by a \emph{domain-altering} expression (\texttt{SUBSTRING}, \texttt{CONCAT}, case mapping) is a new string with no rank and falls back to \cref{piop:lex-sorted-check}; the same holds for comparisons against a literal that need not occur in the column. In TPC-H, nine of the ten queries that order by a string column sort on copies of base columns (Q1, Q2, Q4, Q7, Q9, Q12, Q16, Q20, Q21) and are covered by \cref{piop:rank-sort}. The exception is Q22, whose key $\texttt{cntrycode} = \texttt{SUBSTRING}(\texttt{c\_phone}, 1, 2)$ is a freshly constructed string: \texttt{SUBSTRING} is a domain-altering expression outside the scope of our string protocols, so neither the rank path nor \cref{piop:lex-sorted-check} applies to it directly, as there is no trusted arithmetization of the new column to apply them to. A fixed-length prefix, however, needs no string machinery at all: the two projections $\mathsf{c}_0, \mathsf{c}_1$ of \texttt{c\_phone} (\cref{piop:char-proj}) pack into the single integer $\mathsf{c}_0 \mult 2^{\alpha} + \mathsf{c}_1$, an injective and order-preserving encoding of the key, so the \texttt{GROUP BY}, the \texttt{IN} list and the \texttt{ORDER BY} of Q22 all become integer operations of $\truthtable$, and the verifier decodes the integer back into two characters when opening the result.\ali{Verify the query list against the TPC-H specification, and whether $\truthtable$ dropped these \texttt{ORDER BY} clauses.}

